\chapter{Фермионная спектральная мера оператора Вайнберга на \texorpdfstring{$H_{15}$}{H15}}

Предыдущая глава получила каноническое семейное направление $P_0$ для
майорановского оператора высшей степени. Остаётся проверить более сильное
условие: определяет ли общая мера родителя его физическую амплитуду.

\section{Точная структура литературного механизма}

Фермионное спектральное действие можно разложить на чётную и нечётную
части спектральной функции:
\begin{equation}
 g(z)=f(z^2)+z h(z^2).
 \label{eq:v5-fermionic-even-odd}
\end{equation}
Нечётная часть даёт обычное фермионное действие и нормирует кинетический
член. В пятнадцатимерной геометрии без правого нейтрино чётная
нескалярная часть следующего порядка может дать хиггс-квадратичный
майорановский член.

В конструкции М.~Сакеллариаду и А.~Ситарза нескалярный эндоморфизм на
левом лептонном дублете имеет слабую часть
\begin{equation}
 \tau_{
m w}=i\sigma_2J,
\end{equation}
где $J$ --- вещественная структура. Полученный член пропорционален
\begin{equation}
 (Y_eY_e^*),(LH)(LH)
\end{equation}
и после нарушения электрослабой симметрии даёт майорановскую массу,
зависящую от масштаба отсечения, вакуумного значения Хиггса, масс
заряженных лептонов и коэффициентов функции отсечения; см.
\path{arXiv:1903.09149}.

Статья прямо оставляет открытой классификацию допустимых нескалярных
функций. Для межсемейного смешивания требуется дополнительная
недиагональная семейная часть $\tau$. Поэтому наличие оператора ещё не
равносильно выводу его коэффициента.

\section{Независимость чётного и нечётного моментов}

Текущий родитель не содержит аксиомы, связывающей $f$ и $h$ в
\eqref{eq:v5-fermionic-even-odd}. Это можно проверить явным семейством
быстро убывающих функций
\begin{equation}
 g_\alpha(z)=z e^{-z^2}+\alpha e^{-z^2}.
 \label{eq:v5-fermionic-function-counterfamily}
\end{equation}
Для всех $\alpha$ нечётная часть одинакова:
\begin{equation}
 \frac{g_\alpha(z)-g_\alpha(-z)}2=z e^{-z^2},
\end{equation}
тогда как чётная часть равна
\begin{equation}
 \frac{g_\alpha(z)+g_\alpha(-z)}2=\alpha e^{-z^2}.
\end{equation}
Следовательно, обычная кинетическая нормировка может быть одной и той же,
а коэффициент оператора Вайнберга непрерывно меняться и даже обращаться в
нуль.

\begin{nogocriterion}[Независимость фермионного чётного момента]
Пока родительская мера не связывает чётную и нечётную части одной
фермионной спектральной функции, нормировка обычного действия не определяет
коэффициент оператора Вайнберга. Выбор конкретного $\alpha$ после сравнения
с нейтринной массой является подгонкой спектральной функции.
\end{nogocriterion}

Этот запрет согласуется с ранней заморозкой общей меры: исходный
функционал оставлял свободными функцию отсечения, её масштаб и полную
фермионно-призрачную меру.

\section{Что делает сжатие семейного коэффициента}

Литературный член содержит положительный семейный оператор
\begin{equation}
 K_e=Y_eY_e^*.
\end{equation}
Сжатие на голономную нулевую линию действительно удаляет его матричную
неоднозначность:
\begin{equation}
 P_0K_eP_0=\kappa_eP_0,
 \qquad
 \kappa_e=\langle v_0,K_ev_0\rangle.
 \label{eq:v5-charged-lepton-compression}
\end{equation}
Например, для диагонального
$K_e=\operatorname{diag}(d_1,d_2,d_3)$ получаем
\begin{equation}
 \kappa_e=\frac{d_1+d_2+d_3}{3}.
\end{equation}
Семейное направление снова единственно, но число $\kappa_e$ не выведено:
предыдущие одноформенный и торсионный гейты доказали, что амплитуды
заряженных рёбер $u,d,e$ остаются входными данными.

Поэтому использование наблюдаемых масс заряженных лептонов в
\eqref{eq:v5-charged-lepton-compression} создало бы условную численную
связь, а не предсказание родителя.

\section{Почему вес \texorpdfstring{$1/7$}{1/7} не является массой}

Вес нулевой линии $w_0=1/7$ выведен из следа $M_{35}$. Однако один и тот
же след умножает кинетический и майорановский члены:
\begin{equation}
 S_0=w_0\left[
 Z_\psi\,\overline\psi D\psi
 +\mu\,\psi^TC\psi
 \right].
 \label{eq:v5-common-weight-action}
\end{equation}
После канонической нормировки
\begin{equation}
 \psi_c=\sqrt{w_0Z_\psi}\,\psi
\end{equation}
физическая масса равна
\begin{equation}
 m_{\rm phys}=\frac{\mu}{Z_\psi}
\end{equation}
и от $w_0$ не зависит.

\begin{proposition}[След выбирает канал, но не масштаб]
Число $1/7$ является канонической кратностью голономной нулевой линии в
общем следе $M_{35}$. Оно не может определить нейтринную массу, поскольку
сокращается между кинетическим и массовым членами после нормировки поля.
\end{proposition}

\section{Остаточная формула амплитуды}

С точностью до соглашений теплового разложения коэффициент имеет
структуру
\begin{equation}
 m_{\nu,0}\sim
 r_\tau\,\kappa_e\,\frac{v_H^2}{\Lambda},
 \label{eq:v5-weinberg-residual-amplitude}
\end{equation}
где $r_\tau$ есть отношение нескалярного чётного спектрального момента к
моменту, нормирующему обычную фермионную кинетику. Три множителя остаются
независимыми:
\begin{enumerate}
 \item $r_\tau$ не фиксирован общей мерой;
 \item $\kappa_e$ зависит от невыведенного заряженно-лептонного оператора;
 \item размерный масштаб $\Lambda$ не получен вариационно.
\end{enumerate}
Ни $C_3$, ни проектор $P_0$, ни вес $1/7$ не устраняют эти свободы.

Кроме того, \eqref{eq:v5-weinberg-residual-amplitude} является однородным
эффективным массовым членом. Он не порождает корневую связь, неоднородный
вихрь или радиальную длину локализации.

\begin{theorem}[Незамкнутость фермионной меры]
Нескалярное фермионное спектральное действие даёт допустимый механизм
оператора Вайнберга на $H_{15}$, а голономный проектор фиксирует его
семейное направление. Но текущая родительская мера не связывает чётный
спектральный момент с кинетическим, не выводит заряженно-лептонный
квадратичный коэффициент и не фиксирует масштаб отсечения. Поэтому
параметрически чистая майорановская амплитуда не получена.
\end{theorem}

\section{Критерий остановки}

Дальнейшее варьирование функций отсечения внутри текущего родителя
запрещено: контрсемейство \eqref{eq:v5-fermionic-function-counterfamily}
уже доказывает непрерывную неединственность. Переоткрытие возможно только
после нового принципа меры, который одновременно определяет $f$, $h$,
$\Lambda$ и заряженно-лептонный оператор до обращения к массам.

Следующая глава должна подвести итог части II: отделить полученную
кинематику, канонический семейный канал и условный оператор высшей степени
от невыведенных меры, масштаба и локализации.

\begin{protocol}
\begin{enumerate}
 \item нескалярный оператор на $H_{15}$: \textbf{допустим};
 \item семейное направление $P_0$: \textbf{фиксировано};
 \item вес нулевой линии: \textbf{$1/7$};
 \item влияние веса $1/7$ на физическую массу: \textbf{сокращается};
 \item связь чётного и нечётного спектральных моментов:
 \textbf{не выведена};
 \item коэффициент $\kappa_e$: \textbf{не выведен};
 \item масштаб $\Lambda$: \textbf{не выведен};
 \item дефектная локализация: \textbf{не выведена};
 \item физическое замыкание: \textbf{не достигнуто}.
\end{enumerate}
\end{protocol}

Машинный протокол сохранён в
\path{s2t_v5_h15_fermionic_spectral_weinberg_measure_gate_results.json}.